\documentclass[11pt]{article}
\usepackage{mathblatt}
\begin{document}
\blattkopf*{Basisaufgaben}{BAS-Z6}{Basisaufgaben · Zettel · BAS-Z6}
\einheitenkopf[][Zettel 6]{Basisaufgaben · Zettel 6}
\zweigzeile{ohne Taschenrechner · je Aufgabe 1 Punkt · Lösungen auf der Rückseite}
% \small: zehn Aufgaben mit bis zu zwei Koordinatensystemen auf einer Seite (Render 28.09.)
\small

% 1: Bruchteil einer Fläche bestimmen – brueche-dezimalzahlen-basis-k1-v5 (Original 2026-FOR-B1b)
\begin{aufgabe}{Bei welcher Figur ist genau $\frac{3}{4}$ grau? \hfill (P26F) \\ \kreuz{P} \kreuz{Q} \kreuz{R} \kreuz{S}}

P \bruchkreis[0.8]{3}{5} \quad Q \kreissektor[0.8]{270}{} \quad R \bruchrechteck[2.5]{4}{6} \quad S \bruchrechteck[2.5]{3}{8}
\end{aufgabe}

% 2: Zahlen in verschiedenen Darstellungen vergleichen – brueche-dezimalzahlen-basis-k4-v5 (Original 2023-OS-B1f)
\begin{aufgabe}{Welcher Wert ist der größte: $1{,}2^2$; $1{,}4$; $140\,\%$; $1{,}3$? \hfill (P23) \\ \leerfeld}
\end{aufgabe}

% 3: Termwert berechnen – rationale-zahlen-basis-k1-v4 (Original 2026-FOR-B1g)
\begin{aufgabe}{Berechne den Wert von $3 \cdot (x + 4)$ für $x = -6$. \hfill (P26F) \\ \leerfeld}
\end{aufgabe}

% 4: Exponent einer Potenz bestimmen – potenzen-wurzeln-basis-k1-v2 (Original 2024-OS-B1g)
\begin{aufgabe}{$3^x = 81$ – welche Zahl ist $x$? \hfill (P24) \\ x = \leerfeld}
\end{aufgabe}

% 5: Grundwert berechnen – prozentrechnung-basis-k1-v2 (Original 2025-OS-B1a)
\begin{aufgabe}{Beim Kauf eines Rucksacks spart Ben $9\,€$, das sind $25\,\%$ des alten Preises. Wie hoch war der alte Preis? \hfill (P25) \\ \leerfeld[€]}
\end{aufgabe}

% 6: Prozentwert berechnen – prozentrechnung-basis-k4-v4 (Original 2026-FOR-B1a)
\begin{aufgabe}{$10\,\%$ von $37\,€$ \hfill (P26F) \\ \leerfeld[€]}
\end{aufgabe}

% 7: Term zu Sachtext angeben – terme-basis-k2-v3 (Original 2023-OS-B1h)
\begin{aufgabe}{Das Vierfache einer Zahl $x$ wird um 7 vermindert. Welcher Term passt? \hfill (P23) \\ \kreuz{$4 \cdot (x - 7)$} \kreuz{$7 - 4x$} \kreuz{$4x - 7$} \kreuz{$4x - 7x$}}
\end{aufgabe}

% 8: Rechteckseite aus Umfang berechnen – flaechen-basis-k1-v1 (Original 2018-OS-B1f)
\begin{aufgabe}{Ein Rechteck hat den Umfang $30$ cm und die Seite $a = 9$ cm. Wie lang ist die Seite $b$? \hfill (P18) \\ \leerfeld[cm]}
\end{aufgabe}

% 9: Winkelfunktion Seitenverhältnis angeben – trigonometrie-basis-k2-v3 (Original 2025-OS-B1g)
\begin{aufgabe}{\begin{minipage}[t]{0.6\linewidth}Rechter Winkel bei $C$. Trage den Bruch ein. \hfill (P25) \\ $\mathrm{tan}\,\beta =$ \leerfeld\end{minipage}\hfill\begin{minipage}[t]{0.37\linewidth}\vspace{0pt}
\dreieck{(0,0)}{(3.75,0)}{(1.35,1.8)}{s}{t}{u}{}{\beta}{}
\end{minipage}}
\end{aufgabe}

% 10: Wahrscheinlichkeit einstufig – bruchrechnung-basis-k1-v2 (Original 2014-OS-B1d)
\begin{aufgabe}{Ein Spielwürfel wird einmal geworfen. Wie groß ist die Wahrscheinlichkeit, weder eine 2 noch eine 5 zu werfen? \hfill (P14) \\ \leerfeld}
\end{aufgabe}

\begleitteil
\erg{1}{Q, denn $\frac{270}{360} = \frac{3}{4}$}
\erg{2}{$1{,}2^2 = 1{,}44$; die anderen sind $1{,}4$, $1{,}4$ und $1{,}3$}
\erg{3}{$3 \cdot (-6 + 4) = 3 \cdot (-2) = -6$}
\erg{4}{$3, 9, 27, 81$: vier Faktoren, also $x = 4$}
\erg{5}{$9 \cdot 4 = 36\,€$ (nicht $25\,\%$ von $9\,€$)}
\erg{6}{$0{,}1 \cdot 37 = 3{,}70\,€$}
\erg{7}{$4x - 7$}
\erg{8}{$30 : 2 = 15$; $15 - 9 = 6$ cm}
\erg{9}{$\mathrm{tan}\,\beta = \frac{t}{s}$}
\erg{10}{$\frac{4}{6} = \frac{2}{3}$ (die 1, 3, 4 und 6)}
\end{document}
